\SourceLeaf{150}theile vereinigt geniesst, nemlich die M.\,d.\,kl.\,Q.\ wird in allen Fällen und bei jedem Fehlergesetz die absolut-vortheilhafteste, und die Vergleichung der Genauigkeit der Resultate mit der der Beobachtungen, die ich in meiner Theoria auf das Fehlergesetz $e^{-hhxx}$ gegründet hatte, bleibt allgemein gültig. Zugleich hat man den Vortheil, dass alles durch sehr klare einfache analytische Entwickelungen bewiesen und aufgelöst wird, was bei \textsc{Laplaces} Princip und Behandlung keinesweges der Fall zu sein scheint, so wie namentlich die Generalisirung seines Schlusses von 2 unbekannten Grössen auf jede Anzahl noch nicht die nöthige Evidenz zu haben scheint.

\noindent . . . . . . . . . . .



\Needspace{6\baselineskip}
\begin{center}
\textsc{Gauss} an \textsc{Schumacher}. Göttingen, 2. Februar 1825.
\end{center}

\noindent . . . . . . . . . . Die sogenannten wahrscheinlichen Fehler wünsche ich eigentlich, als von Hypothese abhängig, ganz proscribirt; man mag sie aber berechnen, indem man die mittlern mit 0,6744897 multiplicirt.

\noindent . . . . . . . . . . .



\Needspace{6\baselineskip}
\begin{center}
\textsc{Gauss} an \textsc{Olbers}. Göttingen, 15. März 1827.
\end{center}

\noindent Ich bin Ihnen sehr verpflichtet, mein allertheuerster Freund, für die gefällige Übersendung der drei Hefte des Philosophical Magazine. Ich habe die Aufsätze von \textsc{Ivory} durchgesehen, und ich würde sagen müssen, dass ich dadurch sehr befremdet sei, wenn ich nicht durch die Abhandlung in den Philosophical Transactions (1824), deren ich in meinem letzten Briefe erwähnte, schon vorbereitet gewesen wäre. Ich hatte Herrn I.\ längst als einen scharfsinnigen Mathematiker geschätzt, der im Calcul grosse Gewandtheit hat: von \textls{diesem} Urtheil gehe ich auch noch nicht ab; allein rücksichtlich des Hauptpunkts in der erwähnten Abhandlung vermisse ich den logischen Zusammen%
